%% Ridwan Amure — Academic CV
%% Compile with: pdflatex ridwan-amure-cv.tex

\documentclass[11pt, letterpaper]{article}

\usepackage[margin=1in, top=0.85in, bottom=0.85in]{geometry}
\usepackage[hidelinks]{hyperref}
\usepackage{enumitem}
\usepackage{titlesec}
\usepackage{microtype}
\usepackage{array}
\usepackage{xcolor}
\usepackage{parskip}

%% -------------------------------------------------------------------
%% Typography
%% -------------------------------------------------------------------
\usepackage[T1]{fontenc}
\usepackage{lmodern}

%% -------------------------------------------------------------------
%% Colors
%% -------------------------------------------------------------------
\definecolor{headingblue}{RGB}{30, 64, 110}
\definecolor{mutedgray}{RGB}{100, 100, 100}
\definecolor{rulecolor}{RGB}{180, 180, 180}

%% -------------------------------------------------------------------
%% Section headings
%% -------------------------------------------------------------------
\titleformat{\section}
  {\large\bfseries\color{headingblue}}
  {}{0em}{}
  [\vspace{1pt}\textcolor{rulecolor}{\hrule height 0.6pt}\vspace{4pt}]

\titlespacing{\section}{0pt}{14pt}{6pt}

%% -------------------------------------------------------------------
%% List settings
%% -------------------------------------------------------------------
\setlist[itemize]{
  leftmargin=1.2em,
  itemsep=1pt,
  topsep=3pt,
  parsep=0pt
}

%% -------------------------------------------------------------------
%% Custom commands
%% -------------------------------------------------------------------
\newcommand{\cvheading}[2]{%
  \noindent\textbf{#1} \hfill \textcolor{mutedgray}{\small #2}%
}

\newcommand{\cvsub}[2]{%
  \noindent\textit{#1} \hfill \textcolor{mutedgray}{\small #2}%
}

\newcommand{\pubentry}[1]{%
  \item #1
}

\hypersetup{
  colorlinks=true,
  linkcolor=headingblue,
  urlcolor=headingblue
}

%% -------------------------------------------------------------------
\begin{document}
\pagestyle{empty}

%% ===================================================================
%% HEADER
%% ===================================================================
\begin{center}
  {\LARGE\bfseries Ridwan Amure}\\[6pt]
  {\color{mutedgray}
    Little Rock, Arkansas, United States\\[3pt]
    \href{mailto:amureridwan002@gmail.com}{amureridwan002@gmail.com}
    \enspace\textbar\enspace
    \href{https://ridwanamure.com}{ridwanamure.com}
    \enspace\textbar\enspace
    \href{https://linkedin.com/in/ridwan-amure}{linkedin.com/in/ridwan-amure}\\[2pt]
    \href{https://scholar.google.com/citations?user=pbfF8dQAAAAJ}{Google Scholar}
    \enspace\textbar\enspace
    \href{https://github.com/instabaines}{github.com/instabaines}
    \enspace\textbar\enspace
    \href{https://orcid.org/0000-0001-5874-5871}{ORCID: 0000-0001-5874-5871}
  }
\end{center}

\vspace{4pt}

%% ===================================================================
%% RESEARCH INTERESTS
%% ===================================================================
\section{Research Interests}

Computational Social Science, Network Science, Narrative Diffusion, Graph Mining, Knowledge Graphs, Temporal Modeling, Representation Learning, Large-Scale Machine Learning, Applied AI.

%% ===================================================================
%% EDUCATION
%% ===================================================================
\section{Education}

\cvheading{University of Arkansas at Little Rock}{Little Rock, AR}\\
\cvsub{Ph.D. in Computer Science}{Expected May 2027}
\begin{itemize}
  \item Dissertation: Cross-platform narrative diffusion and collective behavior modeling
  \item Advisor: Prof.\ Nitin Agarwal (Jerry L.\ Maulden-Entergy Endowed Chair, Director of COSMOS)
  \item Research Lab: COSMOS — Collaboratorium for Social Media and Online Behavioral Studies
\end{itemize}

\vspace{6pt}

\cvheading{University of Arkansas at Little Rock}{Little Rock, AR}\\
\cvsub{M.S. in Computer Science}{May 2025}

%% ===================================================================
%% RESEARCH EXPERIENCE
%% ===================================================================
\section{Research Experience}

\cvheading{Graduate Researcher, COSMOS Research Center}{August 2023 -- Present}\\
\cvsub{University of Arkansas at Little Rock}{Little Rock, AR}
\begin{itemize}
  \item Built COSMOS, a large-scale AI pipeline for cross-platform social media analysis, processing millions of posts across datasets with over 100K users.
  \item Developed knowledge graph and embedding-based models for narrative tracking and entity representation across platforms.
  \item Designed ML systems for modeling information diffusion and behavioral dynamics using graph-based and statistical approaches.
  \item Implemented scalable data pipelines for preprocessing, feature extraction, and modeling on high-dimensional datasets.
  \item Applied network science and machine learning to extract actionable insights from complex relational data.
\end{itemize}

%% ===================================================================
%% INDUSTRY EXPERIENCE
%% ===================================================================
\section{Industry Experience}

\cvheading{Meta (Menlo Park, CA)}{Summer 2026}\\
\cvsub{Machine Learning Intern — Ads Ranking}{}
\begin{itemize}
  \item Improving the state-of-the-art early-stage ranking model for Product Ads.
  \item Working on product quantization and target-aware attention for recommendation quality and retrieval efficiency.
  \item Advancing ranking model quality for large-scale ads recommendation systems in production.
\end{itemize}

\vspace{6pt}

\cvheading{Meta (Menlo Park, CA)}{Summer 2025}\\
\cvsub{Software Engineering Intern, Machine Learning}{}
\begin{itemize}
  \item Built end-to-end ML pipeline reducing training data latency from 14 days to 1 day, enabling near real-time model iteration.
  \item Improved large-scale ranking systems via feature engineering, feature ablation, and integration of new high-impact signals.
  \item Designed transformer-based interaction architecture for retrieval and ranking systems, optimizing inference-time performance.
  \item Developed scalable data pipelines supporting training and evaluation of production ML models on large datasets.
\end{itemize}

\vspace{6pt}

\cvheading{IQuartic (NewWave Company)}{January 2023 -- August 2023}\\
\cvsub{Senior Machine Learning Engineer}{Remote}
\begin{itemize}
  \item Developed end-to-end ML systems for ICD-to-HCC risk adjustment, increasing financial gains by over 30\% through improved coding accuracy.
  \item Built Bayesian inference models to identify missed or suboptimal medical coding in electronic health records.
  \item Designed NLP and RAG pipelines for structured extraction and analysis of clinical text data.
\end{itemize}

\vspace{6pt}

\cvheading{ParallelScore}{September 2021 -- April 2023}\\
\cvsub{Data Science Engineer}{Remote}
\begin{itemize}
  \item Developed transformer-based NLP models for medical coding, improving PHI detection from below 50\% to 95\% in QA and 89\% in production.
  \item Designed and deployed inference systems using AWS, Docker, and Kubernetes for real-time model serving.
  \item Conducted model evaluation, error analysis, and performance optimization for production-grade AI systems.
\end{itemize}

\vspace{6pt}

\cvheading{Babban Gona}{September 2020 -- September 2021}\\
\cvsub{Machine Learning Engineer}{Lagos, Nigeria}
\begin{itemize}
  \item Built computer vision models deployed to over 200,000 users for agricultural diagnostics.
  \item Developed low-latency inference systems for real-time detection of plant conditions.
\end{itemize}

%% ===================================================================
%% PUBLICATIONS
%% ===================================================================
\section{Publications}

\textit{* indicates equal contribution. Publications listed in reverse chronological order.}

\vspace{6pt}

\textbf{Journal Articles}

\begin{enumerate}[leftmargin=1.5em, itemsep=5pt, topsep=4pt]

  \pubentry{%
    \textbf{Amure, R.}, \& Agarwal, N. (2026).
    Modeling cross-platform narrative diffusion: a stable Hawkes--ODE framework for exposure and adoption dynamics.
    \textit{Social Network Analysis and Mining}. Springer.
  }

  \pubentry{%
    \textbf{Amure, R.}, \& Agarwal, N. (2026).
    Identifying cohesive narrative formations in cross-platform trade-war discourse using weighted contextual focal structures analysis.
    \textit{Social Network Analysis and Mining}, 16, 49.
    \href{https://doi.org/10.1007/s13278-026-01591-7}{doi:10.1007/s13278-026-01591-7}
  }

  \pubentry{%
    \textbf{Amure, R.}, \& Agarwal, N. (2026).
    Discovering cross-platform narrative flow templates using frequent subgraph mining.
    \textit{Social Network Analysis and Mining}. Springer Vienna.
    \href{https://link.springer.com/content/pdf/10.1007/s13278-025-01577-x_reference.pdf}{[link]}
  }

  \pubentry{%
    \textbf{Amure, R.}, \& Agarwal, N. (2025).
    Co-commenters as clues: a partial-label approach to detecting anomalous channels on YouTube.
    \textit{Social Network Analysis and Mining}, 15(1), 80.
    \href{https://link.springer.com/article/10.1007/s13278-025-01493-0}{[link]}
  }

  \pubentry{%
    \textbf{Amure, R.}, \& Agarwal, N. (2025).
    Modeling cross-platform narrative templates: a temporal knowledge graph approach.
    \textit{Social Network Analysis and Mining}, 15(1), 42.
    \href{https://link.springer.com/article/10.1007/s13278-025-01429-8}{[link]}
  }

  \pubentry{%
    \textbf{Amure, R.}, \& Agarwal, N. (2025).
    A comparative evaluation of social network analysis tools: performance and community engagement perspectives.
    \textit{Social Network Analysis and Mining}, 15(1), 1.
    \href{https://link.springer.com/article/10.1007/s13278-025-01409-y}{[link]}
  }

  \pubentry{%
    Nan, Z., Addai, E., \textbf{Amure, R.}, Bottey, R., Larrey, E.~K., \& Ngungu, M. (2025).
    Memory-driven modeling of herpes simplex virus type-1 and type-2 dynamics with neural network optimization.
    \textit{Computer Methods and Programs in Biomedicine}, 109121. Elsevier.
    \href{https://www.sciencedirect.com/science/article/pii/S0169260725005371}{[link]}
  }

  \pubentry{%
    Akinnubi, A., Alassad, M., \textbf{Amure, R.}, \& Agarwal, N. (2024).
    KG-CFSA: a comprehensive approach for analyzing multi-source heterogeneous social network knowledge graph.
    \textit{Social Network Analysis and Mining}, 14(1), 159.
    \href{https://link.springer.com/article/10.1007/s13278-024-01320-y}{[link]}
  }

  \pubentry{%
    \textbf{Amure, R.}, Alao, O.~A., \& Falebita, D. (2021).
    Lithology identification and prediction using clustering algorithms.
    \textit{NAPE Bulletin}, 30(2), 83--90.
    \href{https://www.researchgate.net/publication/359617732}{[link]}
  }

\end{enumerate}

\vspace{8pt}
\textbf{Conference Papers}

\begin{enumerate}[leftmargin=1.5em, itemsep=5pt, topsep=4pt]

  \pubentry{%
    \textbf{Amure, R.}, \& Agarwal, N. (2026).
    Focal Collective Actors as Narrative Structures: Cross-Platform Brokerage in Digital Protest.
    \textit{Proceedings of the Pacific Asia Conference on Information Systems (PACIS) 2026}.
  }

  \pubentry{%
    \textbf{Amure, R.}, \& Agarwal, N. (2026).
    CI-FSA: Toward Scalable Discovery of Influential Groups in Social Networks.
    In Cherifi et al.\ (Eds.), \textit{Complex Networks \& Their Applications XIV}, pp.\ 108--119. Springer Nature Switzerland.
    \href{https://doi.org/10.1007/978-3-032-16645-6_10}{doi:10.1007/978-3-032-16645-6\_10}
  }

  \pubentry{%
    \textbf{Amure, R.}, \& Agarwal, N. (2026).
    Analysis of Cross-Platform Narrative Dissemination Through Contextual Focal Structures.
    \textit{Social Networks Analysis and Mining (ASONAM) 2025}, 16323, pp.\ 119--127. Springer Nature Switzerland.
    \href{https://link.springer.com/chapter/10.1007/978-3-032-13821-7_11}{[link]}
  }

  \pubentry{%
    \textbf{Amure, R.}, \& Agarwal, N. (2026).
    Exposure, skepticism, and switching: a behavioral model of narrative diffusion.
    \textit{Hawaii International Conference on System Sciences (HICSS) 2026}.
    \href{https://hdl.handle.net/10125/112076}{[link]}
  }

  \pubentry{%
    \textbf{Amure, R.}, \& Agarwal, N. (2025).
    Modeling Word-Level Functions in Social Movement Discourse.
    \textit{International Conference on Information Systems (ICIS) 2025}.
    \href{https://aisel.aisnet.org/icis2025/user_behav/user_behav/30/}{[link]}
  }

  \pubentry{%
    \textbf{Amure, R.}, \& Agarwal, N. (2025).
    Graph-based identification of influential groups in cross-platform narrative dissemination.
    \textit{2025 IEEE International Conference on Knowledge Graph (ICKG)}, pp.\ 9--16. IEEE.
    \href{https://ieeexplore.ieee.org/abstract/document/11391778/}{[link]}
  }

  \pubentry{%
    \textbf{Amure, R.}, \& Agarwal, N. (2025).
    Evaluating Synthetic Data Generation Methods for Anomalous Channel Detection in Sparse-Label Environments.
    \textit{18th International Conference on Social Computing, Behavioral-Cultural Modeling, \& Prediction and Behavior Representation in Modeling and Simulation (SBP-BRiMS) 2025}. SBP-BRiMS.
    \href{https://sbp-brims.org/2025/papers/working-papers/2025_SBP-BRiMS_paper_57.pdf}{[link]}
  }

  \pubentry{%
    \textbf{Amure, R.}, \& Agarwal, N. (2025).
    Modeling Cross-Platform Narrative Diffusion: A Multiplex Approach to Information Spread in Social Media Ecosystems.
    \textit{International Conference on Advances in Social Networks Analysis and Mining (ASONAM) 2025}, pp.\ 166--173. Springer Nature Switzerland.
    \href{https://link.springer.com/chapter/10.1007/978-3-032-13513-1_14}{[link]}
  }

  \pubentry{%
    \textbf{Amure, R.}, \& Agarwal, N. (2024).
    Anomalous channel detection for YouTube through label propagation.
    \textit{International Conference on Complex Networks and Their Applications (CNA) 2024}, pp.\ 271--281. Springer Nature Switzerland.
    \href{https://link.springer.com/chapter/10.1007/978-3-031-82435-7_22}{[link]}
  }

  \pubentry{%
    Shajari, S., \textbf{Amure, R.}, \& Agarwal, N. (2024).
    Detecting and measuring anomalous behaviors on YouTube.
    \textit{International Conference on Advances in Social Networks Analysis and Mining (ASONAM) 2024}, pp.\ 3--16. Springer Nature Switzerland.
    \href{https://link.springer.com/chapter/10.1007/978-3-031-85386-9_1}{[link]}
  }

  \pubentry{%
    Shajari, S., \textbf{Amure, R.}, \& Agarwal, N. (2024).
    Navigating the anomalies: A comprehensive analysis of YouTube channel behavior.
    \textit{International Conference on Advances in Social Networks Analysis and Mining (ASONAM) 2024}, pp.\ 79--87. Springer Nature Switzerland.
    \href{https://link.springer.com/chapter/10.1007/978-3-031-85240-4_8}{[link]}
  }

  \pubentry{%
    \textbf{Amure, R.}, Akinnubi, A., \& Koleoso, O. (2024).
    Toward Improved Clustering for Textual Data.
    \textit{International Conference on Computing and Communication Networks (ICCN) 2024}, pp.\ 405--413. Springer Nature Singapore.
    \href{https://link.springer.com/chapter/10.1007/978-981-97-0892-5_32}{[link]}
  }

  \pubentry{%
    Shajari, S., \textbf{Amure, R.}, \& Agarwal, N. (2024).
    Analyzing anomalous engagement and commenter behavior on YouTube.
    \textit{Thirtieth Americas Conference on Information Systems (AMCIS) 2024}.
    \href{https://aisel.aisnet.org/amcis2024/social_comp/social_comput/6/}{[link]}
  }

  \pubentry{%
    Akinnubi, A., Alassad, M., Agarwal, N., \& \textbf{Amure, R.} (2023).
    Identifying contextualized focal structures in multisource social networks by leveraging knowledge graphs.
    \textit{International Conference on Complex Networks and Their Applications (CNA) 2023}, pp.\ 15--27. Springer Nature Switzerland.
    \href{https://link.springer.com/chapter/10.1007/978-3-031-53472-0_2}{[link]}
  }

\end{enumerate}

%% ===================================================================
%% TALKS AND PRESENTATIONS
%% ===================================================================
\section{Talks and Presentations}

\begin{itemize}

  \item \textbf{Agents at Work: Accelerating and Innovating Business Processes with AI Agents.}\\
  National Society of Black Engineers (NSBE) Annual Convention, Baltimore, MD. March 2026.

  \item \textbf{Building and Deploying Retrieval-Augmented Generation (RAG) Systems for Quality Business Impact.}\\
  National Society of Black Engineers (NSBE) Annual Convention, Chicago, IL. March 2025.

  \item \textbf{Machine Learning Regression and Building Large Language Models from Scratch.}\\
  AI Summer of Code 2025, Lagos, Nigeria (Virtual). September 2025.

  \item \textbf{Mathematical Foundations of AI.}\\
  AI Summer of Code 2024, Lagos, Nigeria (Virtual). February 2024.

  \item \textbf{Scientific Modelling Workshop.}\\
  EAGE Near Surface Geoscience Young Professionals (NGYP), Lagos, Nigeria (Virtual). November 2021.

  \item \textbf{Data Science and Machine Learning for Engineers.}\\
  Society of Petroleum Engineers (SPE) NGYP, Lagos, Nigeria (Virtual). November 2021.

  \item \textbf{Python Virtual Classroom.}\\
  EAGE Near Surface Geoscience Young Professionals (NGYP), Lagos, Nigeria (Virtual). August 2020.

\end{itemize}

%% ===================================================================
%% SOFTWARE AND TOOLS
%% ===================================================================
\section{Software and Open-Source Tools}

\begin{itemize}

  \item \textbf{submine} \hfill \href{https://github.com/instabaines/submine}{github.com/instabaines/submine}\\
  A research-grade Python library for frequent subgraph mining with a unified interface over multiple mining backends and graph formats. Designed for reproducible research workflows.

  \item \textbf{Temporal Eigenvector Centrality} \hfill \href{https://github.com/instabaines/temporal_eigen_vector_centrality}{github.com/instabaines/temporal\_eigen\_vector\_centrality}\\
  A C++ implementation with Python bindings for centrality analysis on temporal and multiplex networks, built for large sparse network analysis.

\end{itemize}

%% ===================================================================
%% AWARDS AND HONORS
%% ===================================================================
\section{Awards, Honors, and Sponsorships}

\begin{itemize}
  \item Linux Foundation PyTorch Conference Sponsorship \hfill 2025
  \item AWS re:Invent Conference Sponsorship \hfill 2024
  \item Acxiom Diversity Scholarship, University of Arkansas at Little Rock \hfill 2024
\end{itemize}

%% ===================================================================
%% SKILLS
%% ===================================================================
\section{Technical Skills}

\begin{itemize}[leftmargin=1.2em, itemsep=3pt]
  \item \textbf{Machine Learning:} PyTorch, TensorFlow, Transformers, Retrieval-Augmented Generation (RAG), Agentic AI, Ranking Systems
  \item \textbf{Data Science:} Python, SQL, Pandas, NumPy, Bayesian Models, Knowledge Graphs, Graph Neural Networks
  \item \textbf{Engineering \& DevOps:} AWS, Docker, Kubernetes, Scalable ML Pipelines, Distributed Training
  \item \textbf{NLP:} Information Extraction, Medical Coding Automation, Language Model Workflows, Entity Recognition
\end{itemize}

%% ===================================================================
%% SERVICE
%% ===================================================================
\section{Certifications}

\begin{itemize}
  \item \textbf{Generative AI with Diffusion Models} --- NVIDIA \hfill April 2025
  \item \textbf{Certified Machine Learning Engineer -- Associate} --- AWS \hfill November 2024
  \item \textbf{Getting Started with Accelerated Computing in CUDA C/C++} --- NVIDIA \hfill November 2023
\end{itemize}

%% ===================================================================
%% SERVICE
%% ===================================================================
\section{Professional Service}

\begin{itemize}
  \item \textbf{Peer Reviewer} --- 30 papers total (17 journal articles, 13 conference papers)
  \begin{itemize}[topsep=2pt]
    \item \textit{Journals (17):} Social Network Analysis and Mining, Springer; The Journal of Supercomputing, Springer; Humanities and Social Sciences Communications, Springer Nature; Journal of King Saud University --- Computer and Information Sciences, Elsevier; Scientific Reports, Nature; Applied Network Science, Springer; Telematics and Informatics, Elsevier; and others
    \item \textit{Conferences (13):} ICIS 2026; AMCIS 2025; IEEE VR 2025; SAIS 2025; HAI 2024; Mensch und Computer 2024; ICWSM 2026; SBP-BRiMS 2025, 2026; HICSS 2027
  \end{itemize}
  \item Member, National Society of Black Engineers (NSBE)
\end{itemize}

\end{document}
